\subsection{Extension of local homomorphisms of groups}

DEFINITION 1. --- Given a topological group G and a group G', a local homomorphism of G into G' is a map f from a neighborhood V of the identity element of G into G' such that, for every pair of points x, y in V with xy ∈ V, one has f(xy) = f(x)f(y).

If G' is a topological group and if f is a continuous map, one says that f is a continuous local homomorphism (or local morphism of topological groups).

If G and G' are topological groups, the notion of local isomorphism from G to G' was defined (TG, III, p. 6, def. 2). A local isomorphism from G to G' is a homeomorphism f from a neighborhood V of the identity element of G onto a neighborhood V' of the identity element of G' such that f and the inverse map of f are local homomorphisms.

PROPOSITION 1. --- Let G and G' be topological groups and let p: G → G' be a group homomorphism. For p to make G a covering of G', it is necessary and sufficient that the restriction of p to a suitable neighborhood of the identity element of G be a local isomorphism from G to G'.

The condition is necessary by proposition 3 of TG, III, p. 6. Conversely, suppose that \(p\) induces a homeomorphism from an open neighborhood V of the identity element \(e\) of G onto a neighborhood V' of the identity element of G'. The map \(p\) is then continuous (TG, III, p. 15, prop. 23) and open (TG, III, p. 16, prop. 24). The image H of \(p\) is a subgroup of G′ containing V, hence open and closed in G′ (TG, III, p. 7, corollary). Let N be the kernel of \(p\) ; we have N∩V = \{e\} ; hence N is discrete (loc. cit., prop. 5). Consequently, \(p\) makes G a covering of H, principal with group N (I, p. 100, cor. 3 of th. 1). Since H is open and closed in G', the G'-space (G, p) is a covering.

PROPOSITION 2. --- Let G be a connected topological group, let G' be a group, and let \(f: V \to G'\) be a local homomorphism of G into G', where V is a connected neighborhood of the identity element of G. Then there exists a connected topological group H, a morphism of topological groups \(p\colon\mathrm{H}\to\mathrm{G}\) such that \((\mathrm{H},p)\) is a covering of G and a group homomorphism \(\varphi\colon\mathrm{H}\to\mathrm{G}'\) such that the set of \(y\in p^{-1}(\mathrm{V})\) satisfying \(f(p(y))=\varphi(y)\) is a neighborhood of the identity element in H.

If G' is a topological group and the map f is continuous, such a homomorphism \(\varphi\) is continuous.

Lemma 1. — Let G be a topological group and let V be a connected neighborhood of the identity element e of G. For every neighborhood U of e in G and every \(x \in V\) , there exists an integer \(n \in N\) and elements \(u_{1}, \ldots, u_{n}\) in U such that \(u_{1} \ldots u_{n} = x\) and \(u_{1} \ldots u_{k} \in V\) for every integer k such that \(1 \leqslant k \leqslant n\) .

We may assume that U is open and contained in V. Let A denote the set of \(x \in V\) satisfying the condition of the lemma. If \(x \in A\) and if \(y \in xU \cap V\) then \(y \in A\) , whence \(AU \cap V \subset A\) ; this shows that A is open in V. Let \(x \in V\) such that \(xU^{-1} \cap A \neq \varnothing\) ; we then have \(x \in AU \cap V\) , therefore \(x \in A\) . Consequently, if \(x \in V\) and \(x \notin A\) , \(xU^{-1} \cap A = \varnothing\) and A is closed in V. Since \(e \in A\) and V is connected, it follows that A = V.

Let us now prove the proposition. Denote by \(j\) the map \(g \mapsto (g, f(g))\) from V to \(\mathbf{G} \times \mathbf{G}'\) and let H be the subgroup of \(\mathbf{G} \times \mathbf{G}'\) generated by \(j(\mathrm{V})\) .

Let U be a neighborhood of \(e\) in G, contained in V. Let \(x \in V\) ; by Lemma 1, there exists \(u_1, \ldots, u_n \in U\) such that \(x = u_1 \ldots u_n\) and such that \(u_1 \ldots u_k \in V\) for every integer \(k\) such that \(1 \leqslant k \leqslant n\) . By induction, we have \(f(u_1 \ldots u_k) = f(u_1) \ldots f(u_k)\) for every integer \(k\) , \(1 \leqslant k \leqslant n\) . In particular, \(j(x)\) belongs to the subgroup generated by \(j(U)\) . It follows that H is generated by \(j(U)\) .

Let \(\mathcal{B}\) be the set of subsets of H of the form \(j(\mathrm{U})\) , where U is a neighborhood of \(e\) in V. Let us show that there exists a unique topology on H, compatible with its group structure, for which \(\mathcal{B}\) is a base of the neighborhood filter of the identity element. For this, it suffices to show that the set \(\mathcal{B}\) satisfies the conditions \((\mathrm{GV}_{\mathrm{I}}^{\prime})\) , \((\mathrm{GV}_{\mathrm{II}}^{\prime})\) and \((\mathrm{GV}_{\mathrm{III}}^{\prime})\) from III, p. 4.

Let U be a neighborhood of e in G, contained in V; there exists a neighborhood \(U'\) of e such that \(U' \cdot U' \subset U\) . For every pair x, y of points of \(U'\) , we have \(xy \in V\) and \(f(xy) = f(x)f(y)\) . Consequently, \(j(\mathrm{U}') \in \mathcal{B}\) and \(j(\mathrm{U}') \cdot j(\mathrm{U}') \subset j(\mathrm{U})\) . This shows that the condition \((\mathrm{GV}_{\mathrm{I}}')\) is satisfied.

The set \(U'' = V \cap U^{-1}\) is then a neighborhood of e in V and, for \(x \in U''\) , we have \(x^{-1} \in U\) and \(f(x^{-1}) = f(x)^{-1}\) . Consequently, \(j(\mathrm{U}'')^{-1} \subset j(\mathrm{U})\) , which shows the condition \((\mathrm{GV}_{\mathrm{II}}' )\) .

Finally, let us fix a neighborhood U of e in V such that \(U = U^{-1}\) and \(U^{3} \subset V\) . Let W be a neighborhood of e in U and let \(h = (g, f(g))\) be an element of \(j(\mathrm{U})\) . There exists a neighborhood \(W'\) of e contained in W such that \(gW'g^{-1} \subset W\) . Then, \(j(\mathrm{W}') \in \mathcal{B}\) and \(hj(\mathrm{W}')h^{-1} \subset j(\mathrm{W})\) , since we have \(f(gxg^{-1}) = f(g)f(x)f(g^{-1})\) , for \(x \in W'\) .

Let \(h \in H\) . Since U is a neighborhood of e contained in V, \(j(\mathrm{U})\) generates H; since U is symmetric, there exist elements \(u_{1}, \ldots, u_{n}\) in U such that \(h = j(u_{1}) \ldots j(u_{n})\) . By induction on n, there exists a neighborhood \(W'\) of e contained in W such that \(hj(\mathrm{W}')h^{-1} \subset j(\mathrm{W})\) . Consequently, the condition \((\mathrm{GV}_{\mathrm{III}}')\) is satisfied.

We then endow the group H with this topology.

Denote by \(p \colon \mathrm{H} \to \mathrm{G}\) the restriction to H of the first projection \(\mathrm{G} \times \mathrm{G}' \to \mathrm{G}\) . This is a group homomorphism. For every neighborhood U of \(e\) contained in V, the set \(p^{-1}(\mathrm{U})\) contains the neighborhood \(j(\mathrm{U})\) of the identity element of H, hence \(p\) is a continuous homomorphism of topological groups. For every neighborhood U of \(e\) contained in V, we have \(p(j(\mathrm{U})) = \mathrm{U}\) , therefore \(p\) is an open mapping. Since G is connected, \(p\) is surjective. Its kernel is discrete because it does not meet \(j(\mathrm{V})\) except at the identity element. It then follows from Corollary 3 (I, p. 100) that the G-space \((\mathrm{H}, p)\) is a covering.

Let \(\varphi\) be the restriction to H of the second projection \(G \times G' \to G'\) . If \(g \in V\) , we have \((g, f(g)) = j(g) \in j(V)\) and \(\varphi(g, f(g)) = f(g)\) , so that \(\varphi(y) = f(p(y))\) for \(y \in j(V)\) .

Suppose further that G' is a topological group and that the map f is continuous at e. Then the homomorphism \(\varphi\) is continuous at the identity element of H, hence is continuous (TG, III, p. 15, prop. 23).

COROLLARY 1. --- Let G be a simply connected topological group and let G' be a group. Let V be a connected neighborhood of the identity element in G and let f: V → G' be a local homomorphism. There exists a unique group homomorphism h: G → G' extending f. If G' is a topological group and if the map f is continuous, then the homomorphism h is continuous.

The group G is connected. Let H, \(\varphi\), and p be as in Proposition 2. Since G is simply connected, H is a trivializable covering of G (I, p. 124, def. 3); since H is connected and nonempty, the map \(p\) is an isomorphism of topological groups. Set \(h = \varphi \circ p^{-1}\) ; the map \(h\) is a group homomorphism and there exists an open neighborhood U of the identity element \(e\) of G contained in V such that \(f|U = h|U\) . It follows from Lemma 1 that \(f\) and \(h\) coincide on V. In other words, the map \(h\) extends the map \(f\) .

If G' is a topological group and the map f is continuous, the homomorphism \(\varphi\) is continuous, so h is too.

Let us prove the uniqueness of such an extension. The set of points of G where two homomorphisms from G to G′ coincide is a subgroup of G. Since the group G is connected, every subgroup of G containing a neighborhood of the identity element is equal to G (TG, III, p. 8, prop. 6). The uniqueness of h follows from this.

Remark 1. --- When G = R, this corollary follows from Proposition 6 of TG, V, p. 3.

COROLLARY 2. --- Two locally connected and simply connected topological groups that are locally isomorphic are isomorphic.

Let G and G' be locally connected topological groups, let V and V' be connected neighborhoods of the identity element of G and of G', respectively, and let \(f: V \to V'\) be a homeomorphism that is a local isomorphism from G to G'. By Corollary 1, there exists a unique continuous group homomorphism \(\varphi: G \to G'\) which extends f and a unique continuous group homomorphism \(\varphi': G' \to G\) which extends \(f^{-1}\) . The homomorphisms \(\varphi' \circ \varphi\) and \(\operatorname{Id}_{\mathrm{G}}\) coincide on a neighborhood of e in G, hence are equal, since G is connected. Likewise, \(\varphi \circ \varphi' = \operatorname{Id}_{\mathrm{G}'}\) , whence the corollary.

COROLLARY 3. --- Let G be a simply connected topological group and let V be a connected neighborhood of the identity element of G. One defines a presentation of G by taking as generating set the set V and as the set \(\mathbf{r}\) of relators the family of \(xyz^{-1}\) , where \((x,y,z)\) traverses the triples of elements of V such that \(xy = z\) .

Let \(\mathrm{F}(\mathrm{V}, \mathbf{r})\) be the quotient group of the free group \(\mathrm{F}(\mathrm{V})\) by the smallest normal subgroup containing the elements \(xyz^{-1}\) , where \((x, y, z) \in \mathrm{V} \times \mathrm{V} \times \mathrm{V}\) and xy = z (A, I, p. 86). Let us denote by \(f: V \to F(V, r)\) and \(g: F(V, r) \to G\) the canonical maps. By construction, the map f is a local homomorphism from G to \(\mathrm{F}(V, \mathbf{r})\) . By Corollary 1, there exists a unique group homomorphism \(\overline{f}: G \to F(V, r)\) extending f.

Since the group \(F(V, \mathbf{r})\) is generated by \(f(V)\) , the homomorphism \(\overline{f}\) is surjective. For \(x \in V\) , we have \(g(\overline{f}(x)) = g(f(x)) = x\) ; since V generates G, \(g \circ \overline{f} = \operatorname{Id}_{\operatorname{G}}\) , which proves that \(\overline{f}\) is injective. It is therefore an isomorphism.
% semantic-fragments: legacy-input-seam
\relax{}
